% Source: https://artofproblemsolving.com/community/c3h2670574_some_actually_good_proving_problems
% Stable identity: aops-topic-2670574-q07
\emph{Right triangles.} A right triangle with legs $p$ and $q$ is half of a $p\times q$ rectangle: a diagonal cuts the rectangle into two triangles that are congruent (SSS), so each has area $\dfrac12pq$.

\emph{General triangle.} Let the side be $BC=b$, and let $H$ be the foot of the perpendicular from $A$ to line $BC$, so $AH=h$.

\begin{asy}
usepackage("amsmath");
size(6cm);
defaultpen(fontsize(10pt));
pair B=(0,0),C=(3,0),A=(1.1,2),H=(1.1,0);
draw(A--B--C--cycle,linewidth(0.8));draw(A--H,dashed);draw(H+(0.18,0)--H+(0.18,0.18)--H+(0,0.18));
label("$A$",A,N);label("$B$",B,SW);label("$C$",C,SE);label("$H$",H,S);label("$h$",(A+H)/2,E);
label("Case 1",(1.5,-0.7),fontsize(8pt));
pair B2=(4.2,0),C2=(5.7,0),A2=(7,2),H2=(7,0);
draw(A2--B2--C2--cycle,linewidth(0.8));draw(C2--H2,dashed);draw(A2--H2,dashed);draw(H2+(-0.18,0)--H2+(-0.18,0.18)--H2+(0,0.18));
label("$A$",A2,N);label("$B$",B2,SW);label("$C$",C2,S);label("$H$",H2,SE);label("$h$",(A2+H2)/2,E);
label("Case 2",(5.6,-0.7),fontsize(8pt));
\end{asy}

\emph{Case 1: $H$ lies on segment $BC$.} Then $AH$ cuts the triangle into right triangles $ABH$ and $AHC$ (one of them is absent if $H$ is $B$ or $C$), so
\[[ABC]=\dfrac12\,BH\cdot h+\dfrac12\,HC\cdot h=\dfrac12(BH+HC)h=\dfrac12bh.\]

\emph{Case 2: $H$ lies outside segment $BC$}, say with $C$ between $B$ and $H$. Then triangle $ABH$ is made of the non-overlapping triangles $ABC$ and $ACH$, so
\[[ABC]=[ABH]-[ACH]=\dfrac12\,BH\cdot h-\dfrac12\,CH\cdot h=\dfrac12(BH-CH)h=\dfrac12bh.\ \blacksquare\]
